\section{0. Cómo usar este machote}\label{cuxf3mo-usar-este-machote}

Este documento es la plantilla del \textbf{Documento} \textbf{de}
\textbf{Ingeniería} \textbf{del} \textbf{Proyecto} \textbf{--}
\textbf{Fase} \textbf{1} que cada equipo entrega al profesor para su
evaluación. Reúne en un solo archivo lo que el equipo produce durante
los sprints 0 a 6 y que se encuentra en la carpeta docs/ del
repositorio.

\subsection{0.1 Reglas de uso}\label{reglas-de-uso}

\begin{itemize}
\item
  Conserven la numeración y el orden de las secciones. Si una sección no
  aplica, déjenla con la leyenda "No aplica" y una justificación de una
  línea.
\item
  Los recuadros grises marcados como \textbf{"Instrucción"} explican qué
  debe contener cada sección. \textbf{Bórrenlos antes de entregar.}
\item
  El texto gris entre corchetes, como {[}Escriba aquí\ldots{]}, es un
  marcador: reemplácenlo por contenido propio.
\item
  Las filas de ejemplo en las tablas solo muestran el formato;
  sustitúyanlas por los datos de su proyecto.
\item
  Los diagramas se insertan como imágenes legibles y su archivo fuente
  (draw.io, PlantUML o Mermaid) queda en el repositorio. Cada diagrama
  lleva número, título y ruta de su fuente: \emph{Figura 4.2. Máquina de
  estados de Pedido. Fuente: docs/04-analisis/estados/pedido.puml}.
\item
  Cuando el contenido sea muy extenso (por ejemplo, las 35 historias de
  usuario o los 25 casos de prueba), el documento incluye la tabla
  resumen y el detalle completo puede quedar en un anexo o en el
  repositorio, con su enlace.
\end{itemize}

\subsection{0.2 Entrega acumulativa}\label{entrega-acumulativa}

El documento \textbf{crece sprint a sprint}. En cada Sprint Review se
entrega la versión vigente con las secciones correspondientes completas;
el profesor evalúa solo las secciones del corte, pero puede pedir
corregir las anteriores.

{\def\LTcaptype{none} % do not increment counter
\begin{longtable}[]{@{}
  >{\raggedright\arraybackslash}p{(\linewidth - 6\tabcolsep) * \real{0.1923}}
  >{\raggedright\arraybackslash}p{(\linewidth - 6\tabcolsep) * \real{0.2031}}
  >{\raggedright\arraybackslash}p{(\linewidth - 6\tabcolsep) * \real{0.3909}}
  >{\raggedright\arraybackslash}p{(\linewidth - 6\tabcolsep) * \real{0.2137}}@{}}
\toprule\noalign{}
\begin{minipage}[b]{\linewidth}\raggedright
\textbf{Corte}
\end{minipage} & \begin{minipage}[b]{\linewidth}\raggedright
\textbf{Versión del documento}
\end{minipage} & \begin{minipage}[b]{\linewidth}\raggedright
\textbf{Secciones que deben estar completas}
\end{minipage} & \begin{minipage}[b]{\linewidth}\raggedright
\textbf{Evaluación en la que cuenta}
\end{minipage} \\
\midrule\noalign{}
\endhead
\bottomrule\noalign{}
\endlastfoot
Review 0 (sem. 5) & v0.1 & Portada, 1, 2.1-- 2.4, 3 & Parcial 1 \\
Review 1 (sem. 7) & v0.2 & + 4.1-- 4.5 (borrador), 4.6, 4.7 & Parcial
1 \\
Review 2 (sem. 9) & v1.0 (SRS firmado) & + 4 completa, 5.1--5.3, 2.5 &
Parcial 2 \\
Review 3 (sem. 11) & v1.1 & + 6 & Parcial 3 \\
Review 4 (sem. 13) & v1.2 & + 7.1--7.5, 9.1--9.3 & Parcial 3 \\
Review 5 (sem. 15) & v1.3 & + 7.6--7.8, 8, 9.4 & Ordinaria \\
Demo final (sem. 16) & v2.0 (final de la Fase 1) & Todo el documento,
incluidas 10, 11, 12, 13 y anexos & Ordinaria \\
\end{longtable}
}

\subsection{0.3 Formato de entrega}\label{formato-de-entrega}

\begin{itemize}
\item
  \textbf{Archivos:} este documento en Word (.docx) y en PDF, en
  docs/entregas/ del repositorio, y la liga del PDF en la minuta de la
  Review.
\item
  \textbf{Nombre del archivo:}
  IS1-Fase1-Equipo\textless A\textbar B\textgreater-v\textless versión\textgreater.pdf,
  por ejemplo IS1-Fase1-EquipoA-v1.0.pdf.
\item
  \textbf{Formato:} carta, márgenes de 2.5 cm, fuente de 11 puntos,
  encabezados numerados, numeración de páginas, figuras y tablas
  numeradas.
\item
  \textbf{Citas y referencias:} formato APA 7. Toda norma, libro o dato
  de terceros se cita; las referencias deben ser verificables.
\item
  \textbf{Uso de IA:} se declara en la sección 2.6. Cualquier integrante
  debe poder explicar cualquier parte del documento.
\end{itemize}

\begin{quote}
Si la versión del documento no coincide con lo que hay en el repositorio
y en el tablero, el profesor toma como válido el repositorio.
\end{quote}
